\section{Preliminaries}
% \subsection{Notation}
% \label{sec:notation}
\textbf{Notation.} We use $\langle\cdot,\cdot\rangle$ and $\|\cdot\|$ for the Euclidean inner
product and norm.  For a matrix, $\|\cdot\|$ denotes the spectral norm.  For
$\calA=(A_1,\ldots,A_m)^\top$, let $J_{\calA}(\bx)\in\RR^{m\times n}$ denote its
Jacobian, whose $i$th row is $\nabla A_i(\bx)^\top$.  We write
$\sigma_{\min}(J_{\calA}(\bx))$ for the smallest singular value of
$J_{\calA}(\bx)$. For a nonempty closed convex set $C$, we denote by
$\Pi_C(u):=\argmin_{v\in C}\|u-v\|^2$ the Euclidean projection of $u$ onto $C$.
% Throughout the paper, $\partial h$ denotes the limiting
% subdifferential of $h$.

For a proper closed convex function \(h\), its convex subdifferential and
proximal mapping are, respectively,
\[
    \partial h(y)
    :=\{v:\ h(z)\ge h(y)+\langle v,z-y\rangle\ \text{for all }z\},
\]
\begin{equation}
    \prox_{\gamma h}(u)
    :=\argmin_{y\in\R^m}
       \left\{h(y)+\frac{1}{2\gamma}\|y-u\|^2\right\},
    \qquad \gamma>0.
    \label{eq:prox-definition}
\end{equation}


\subsection{Assumptions}
\label{sec:assumptions}
We impose the following assumptions.
% \begin{assumption}[Level-bounded objective]
% \label{ass:level-boundedness}
% Assume that $F$ is level-bounded, that is, for every $\alpha\in\RR$, the set
% \[
%     \calS_{\alpha}:=\{\bx\in\RR^n:F(\bx)\le\alpha\}
% \]
% is empty or compact.
% \end{assumption}
% \begin{assumption}[Uniform Jacobian regularity]
% \label{ass:strong-licq}
% There exists a constant $\nu>0$ such that
% \[
%     \sigma_{\min}\bigl(J_{\calA}(\bx)\bigr)\ge \nu,
%     \qquad
%     \forall \bx\in\calS.
% \]
% Equivalently,
% \[
%     \|J_{\calA}(\bx)^\top \by\|
%     \ge
%     \nu \|\by\|,
%     \qquad \forall \bx\in\calS,\ \forall \by\in\RR^m.
% \]
% \end{assumption}
\begin{assumption}
\label{ass:smoothness}
There exist constants $L_f,G_f,L_{\calA},G_{\calA}\ge0$ such that, for
all $\bx,\bx'\in\RR^n$,
\[
\begin{aligned}
    \|\nabla f(\bx)-\nabla f(\bx')\|
    &\le L_f\|\bx-\bx'\|,
    &\|\nabla f(\bx)\|&\le G_f,\\
    \|J_{\calA}(\bx)-J_{\calA}(\bx')\|
    &\le L_{\calA}\|\bx-\bx'\|,
    &\|J_{\calA}(\bx)\|&\le G_{\calA}.
\end{aligned}
\]
\end{assumption}

\begin{assumption}
\label{ass:h}
The function $h$ is proper, closed, convex, and globally $L_h$-Lipschitz:
\[
    |h(\bu)-h(\bv)|\le L_h\|\bu-\bv\|,
    \qquad \bu,\bv\in\RR^m.
\]
\end{assumption}

\subsection{Setup}
To study Question 1, we introduce the following problem class.
\begin{definition}[Problem class]
\label{def:problem-class}
Let \(\Delta>0\) and
\(\mathbf C=(L_f,G_f,L_{\calA},G_{\calA},L_h)\in\RR^5_{\ge0}\).
We denote by \(\mathcal I(\mathbf C,\Delta)\) the set of all triples
\((f,\calA,h)\), over all dimensions \(n,m\in\NN\), such that
Assumptions~\ref{ass:smoothness} and~\ref{ass:h} hold with the constants
in \(\mathbf C\) and $F(\zero)-\inf_{\bx\in\RR^n}F(\bx)\le\Delta$, where the constants \(\mathbf C\) and \(\Delta\) is independent on the target accuracy or on the dimensions \(n,m\). 
\end{definition}
% The gap condition implies
% \(\inf F>-\infty\). Assumption~\ref{ass:level-boundedness} is not part of
% \(\mathcal I(\mathbf C,\Delta)\); it is imposed separately in the complexity
% results that require it.

\red{We use the following stationarity measure for  Problem~\eqref{prob:main} as below. }
\begin{definition}[\(\epsilon\)-KKT point]
\label{def:eps-kkt}
Given \(\epsilon>0\). A triple \((\bx,\by,\bq)\) with
\(\bq\in\partial h(\by)\) is an \(\epsilon\)-KKT point of
\eqref{prob:main} if
\begin{equation}
    \mathcal{R}_{\mathrm{KKT}}(\bx,\by,\bq)
    :=
    \max\Bigl\{
        \bigl\|\nabla f(\bx)+J_{\calA}(\bx)^\top\bq\bigr\|,\;
        \|\calA(\bx)-\by\|
    \Bigr\} \leq \epsilon
    \label{eq:kkt-residual}
\end{equation}
\end{definition}

We next relate Definition~\ref{def:eps-kkt} to game stationarity (GS) for the
minimax reformulation~Problem~\eqref{eq:ncl} of Problem~\eqref{prob:main}. We use the
following GS definition from \cite[Definition~3.1(b)]{li2025nonsmooth} and
\cite[(2.10)--(2.11)]{zhang2026matching}.
\begin{definition}[\(\epsilon\)-GS point]
\label{def:eps-gs}
Let \(\epsilon>0\). A pair \((\bx,\bq)\in\RR^n\times\dom\partial h^*\) is an
\(\epsilon\)-game stationary point (\(\epsilon\)-GS) of~Problem~\eqref{eq:ncl} if
\[
    \calG(\bx,\bq)
    :=
    \max\Bigl\{
        \bigl\|\nabla f(\bx)+J_{\calA}(\bx)^\top\bq\bigr\|,\;
        \dist\bigl(\calA(\bx),\partial h^*(\bq)\bigr)
    \Bigr\}
    \le\epsilon .
\]
\end{definition}
We now relate $\epsilon$-KKT points of Problem~\eqref{prob:main}
to $\epsilon$-GS points of its Fenchel saddle reformulation~Problem~\eqref{eq:ncl}. 
\begin{proposition}[KKT points of Problem~\eqref{prob:main} and GS points of~Problem~\eqref{eq:ncl}]
\label{prop:kkt-gs}
Suppose that Assumption~\ref{ass:h} holds, and let
\((\bx,\bq)\in\RR^n\times\dom\partial h^*\). Then the following measurement relationship holds: 
\[
    \calG(\bx,\bq)
    =
    \min_{\by\in\partial h^*(\bq)}
    \mathcal{R}_{\mathrm{KKT}}(\bx,\by,\bq),
\]
with minimum attained at
\(\hat\by=\Pi_{\partial h^*(\bq)}\bigl(\calA(\bx)\bigr)\). Consequently, for
every \(\epsilon>0\):
\begin{enumerate}[label=(\alph*)]
\item if \((\bx,\by,\bq)\) is an \(\epsilon\)-KKT point
  of~Problem~\eqref{prob:main}, then \((\bx,\bq)\) is an \(\epsilon\)-GS point
  of~Problem~\eqref{eq:ncl};
\item if \((\bx,\bq)\) is an \(\epsilon\)-GS point of~Problem~\eqref{eq:ncl}, then
  \((\bx,\hat\by,\bq)\) is an \(\epsilon\)-KKT point of~Problem~\eqref{prob:main}.
\end{enumerate}
\end{proposition}

\begin{proof}
Since \(h\) is proper, closed, and convex,
\(\bq\in\partial h(\by)\) if and only if \(\by\in\partial h^*(\bq)\)
\cite[Theorem~23.5]{rockafellar1970convex}. The first term of
\(\mathcal{R}_{\mathrm{KKT}}\) does not depend on \(\by\). The set
\(\partial h^*(\bq)\) is nonempty, closed, and convex, so
\(\min_{\by\in\partial h^*(\bq)}\|\calA(\bx)-\by\|
=\dist(\calA(\bx),\partial h^*(\bq))\), attained at \(\hat\by\).
Statements~(a) and~(b) follow.
\end{proof}

We next introduce the following deterministic first-order oracle model.
\begin{definition}[Oracle model]
\label{def:oracle-model}
We assume access to a deterministic first-order oracle
\[
    \mathcal{O}_{f,\mathcal{A}}(\bx)
    :=
    \bigl(f(\bx),\nabla f(\bx),\mathcal{A}(\bx),J_{\mathcal{A}}(\bx)\bigr),
\]
and an exact proximal oracle
\[
    \mathcal{O}_{h}(\bv,t)
    :=
    \operatorname{prox}_{t h}(\bv),
    \qquad t>0.
\]
Each call to either oracle counts as one oracle evaluation. In particular,
the cost of solving the proximal subproblem is represented by one exact
proximal-oracle call.
\end{definition}

\subsection{Dual regularization and effective penalties}
\label{sec:dual-penalty-identity}

This subsection clarifies the relationship between Fenchel duality,
augmented Lagrangian regularization, and the dual damping used in our
method. While , we clarify for completeness. 

\paragraph{ALM as a dual proximal update}
For any $\sigma>0$ and $\mu\in\mathbb R^m$, define
\[
\widehat{\mathcal L}_{\sigma}(x,\mu)
:=
f(x)+\min_y
\left\{
h(y)+\langle \mu,A(x)-y\rangle
+\frac{\sigma}{2}\|A(x)-y\|^2
\right\}.
\]
By Fenchel duality,
\[
\widehat{\mathcal L}_{\sigma}(x,\mu)
=
f(x)+\max_{\lambda}
\left\{
\langle A(x),\lambda\rangle-h^*(\lambda)
-\frac{1}{2\sigma}\|\lambda-\mu\|^2
\right\}.
\]
Thus, the standard ALM update with penalty $\rho$ is equivalent to a
dual proximal update with proximal parameter $\rho$ and center
$\mu=\lambda^k$:
\[
\begin{aligned}
\lambda^{k+1}
&=
\operatorname{prox}_{\rho h^*}
\bigl(\lambda^k+\rho A(x^{k+1})\bigr) \\
&=
\argmin_{q\in\mathbb R^m}
\left\{
h^*(q)
+\frac{1}{2\rho}
\left\|
q-\bigl(\lambda^k+\rho A(x^{k+1})\bigr)
\right\|^2
\right\}.
\end{aligned}
\]

\paragraph{Zero center and quadratic penalty}
When the dual center is fixed at $\mu=0$, the same identity reduces to
\[
\max_{\lambda}
\left\{
\langle A(x),\lambda\rangle-h^*(\lambda)
-\frac{1}{2\sigma}\|\lambda\|^2
\right\}
=
h_{1/\sigma}(A(x)),
\]
where $h_{1/\sigma}$ denotes the Moreau envelope of $h$. Hence, the
zero-center formulation is equivalent to the quadratic-penalty value
function
\[
f(x)+h_{1/\sigma}(A(x)).
\]
Its feasibility control is obtained through the penalty parameter
$\sigma$.

\paragraph{Additional dual regularization in our method.}
Our method starts from the reduced augmented Lagrangian
$\mathcal D_\rho(x,\lambda)$ and adds a moving-center dual regularizer:
\[
\mathcal P_s(x,\lambda;z,\mu)
=
\mathcal D_\rho(x,\lambda)
+\frac{1}{2\gamma}\|x-z\|^2
-\frac{1}{2s}\|\lambda-\mu\|^2.
\]
Therefore, the additional quadratic term is applied to the already
reduced augmented Lagrangian, rather than directly to the original
Lagrangian. Maximizing over $\lambda$ and applying the preceding
identity gives
\[
\max_{\lambda}\mathcal P_s(x,\lambda;z,\mu)
=
\widehat{\mathcal L}_{\rho+s}(x,\mu)
+\frac{1}{2\gamma}\|x-z\|^2.
\]
Thus, the dual regularization produces the effective penalty
\[
\rho_{\mathrm{eff}}=\rho+s
\]
while the base ALM penalty $\rho$ remains fixed.

\paragraph{Feasibility interpretation.}
At the exact effective saddle point, let
$q_\sigma\in\partial h(y_\sigma)$ with $\sigma=\rho+s$. Then
\[
A(x)-y_\sigma
=
\frac{q_\sigma-\mu}{\rho+s}.
\]
In particular, when $\mu=0$ and $h$ is $L_h$-Lipschitz,
\[
\|A(x)-y_\sigma\|
\le
\frac{L_h}{\rho+s}.
\]
For a general moving center, the corresponding bound is
\[
\|A(x)-y_\sigma\|
\le
\frac{L_h+\|\mu\|}{\rho+s}.
\]
These are upper bounds induced by the effective penalty. Feasibility can
also be recovered with a bounded $\rho+s$ if the center update drives
$q_\sigma-\mu$ to zero. For the actual single-loop method, this
property must be established from the coupled multiplier and center
updates; the value-function identity alone does not imply exact ALM
equivalence.
